\section{引言：为什么需要脑机接口}

脑机接口（Brain-Computer Interface, BCI）是一项旨在帮助因神经系统疾病而丧失语言或运动能力的患者恢复沟通能力的前沿技术。本次讲座由斯坦福大学神经假体转化实验室（Neural Prosthetics Translational Laboratory, NPTL）的 Chaofei Fan 主讲，系统介绍了语音脑机接口的发展历程、核心技术原理和最新研究成果。

\begin{figure}[H]
\centering
\includegraphics[width=0.95\textwidth]{slides-images/slide-000.jpg}
\caption{课程标题页：Speech Brain-Computer Interfaces for Restoring Natural Communication\protect\footnotemark}
\end{figure}
\footnotetext{Slides 第1页。讲者来自 Stanford NPTL 实验室和 BrainGate 研究联盟。}

\subsection{患者的困境}

许多神经系统疾病——如脑干中风（brainstem stroke）和肌萎缩侧索硬化症（Amyotrophic Lateral Sclerosis, ALS）——可以导致严重的言语和运动障碍，甚至完全丧失语言能力。

\begin{figure}[H]
\centering
\includegraphics[width=0.95\textwidth]{slides-images/slide-002.jpg}
\caption{神经系统疾病患者面临的困境：Howard Wicks 因严重中风失去了所有的梦想\protect\footnotemark}
\end{figure}
\footnotetext{Slides 第3页。}

讲者以 Howard Wicks 为例：这位 21 岁的年轻人因严重中风陷入了``闭锁状态''（locked-in state）——大脑功能完好，却无法移动身体、无法开口说话。对于这类患者来说，与外界沟通是最迫切的需求。

\begin{importantbox}{闭锁状态的核心矛盾}
患者的大脑仍然完全正常运作，拥有完整的思维能力，但身体失去了执行大脑命令的能力。大脑与身体之间的通路被疾病切断了——这正是脑机接口试图``绕过''的环节。
\end{importantbox}

\subsection{现有辅助沟通手段的局限}

目前患者可用的辅助沟通设备主要有两类：

\begin{figure}[H]
\centering
\includegraphics[width=0.95\textwidth]{slides-images/slide-003.jpg}
\caption{辅助沟通设备：字母板（上）和眼动追踪虚拟键盘（下）\protect\footnotemark}
\end{figure}
\footnotetext{Slides 第4页。}

\begin{itemize}
    \item \textbf{字母板}：患者通过残余的眼球运动注视字母，由照护者逐字母拼读。这种方式极其缓慢——说一句简单的话可能需要几分钟。
    \item \textbf{眼动追踪设备}：患者通过注视屏幕上的虚拟键盘打字。虽然比字母板快，但长时间盯着屏幕对患者非常疲劳，且许多患者的残余眼球运动能力也十分有限。
\end{itemize}

\subsection{脑机接口带来的希望}

\begin{figure}[H]
\centering
\includegraphics[width=0.95\textwidth]{slides-images/slide-005.jpg}
\caption{Neuralink 的植入式脑机接口设备及其参与者 Noland Arbaugh 的评语\protect\footnotemark}
\end{figure}
\footnotetext{Slides 第6页。Noland Arbaugh 是 Neuralink PRIME 研究的参与者。}

脑机接口的核心理念是：既然患者的大脑仍然功能完好，我们是否可以\textbf{直接从大脑中读取信号}，绕过损伤的身体通路，让患者重新控制计算机、机械臂，甚至直接恢复语言能力？

\begin{knowledgebox}{脑机接口的定义}
\textbf{脑机接口（BCI）}是一种在大脑与外部设备之间建立直接通信通道的系统。它通过采集和解码大脑神经信号，将用户的意图转化为对外部设备（计算机、机械臂、语音合成器等）的控制命令。BCI 不依赖于身体的自然输出通路（如肌肉运动），因此可以帮助完全瘫痪的患者恢复与外界的交互能力。
\end{knowledgebox}

\subsection{本章小结}

神经系统疾病可以导致患者完全丧失语言和运动能力，但大脑本身往往仍然功能完好。现有辅助沟通手段速度慢、效率低，远不能满足患者的沟通需求。脑机接口提供了一种全新的解决方案——直接从大脑读取信号，为患者搭建通向外部世界的桥梁。

%% ========== Part 2: BCI 简史与神经科学基础 ==========

